\documentclass[11pt]{article}

\usepackage[margin=1in]{geometry}
\usepackage{amsmath}
\usepackage{amssymb}
\usepackage{amsthm}
\usepackage{mathrsfs}
\usepackage{mathtools}
\usepackage{tikz}
\usepackage[colorlinks=true, linkcolor=blue, bookmarks=true]{hyperref}

\renewcommand{\qedsymbol}{$\blacksquare$}

\title{Comprehensive Practice Problem Set:\\ Mathematical Proof Techniques \& Probability Modeling}
\author{Tashfeen Omran}
\date{September 2026}

\begin{document}

\maketitle

\tableofcontents
\newpage

\section{Concept Checklist \& Question Mapping}
\label{sec:checklist_start}

Based on the provided theoretical document, the following logical structures, theorems, probability models, and proof techniques have been extracted and mapped to the generated problem set.

\begin{itemize}
    \item \textbf{C1: Direct Proof (Parity \& Algebra)} -- Problems: 1, 27
    \item \textbf{C2: Proof by Contrapositive (Implications \& De Morgan's)} -- Problems: 2, 6, 20, 26
    \item \textbf{C3: Proof by Cases (Exhaustion \& Partitions)} -- Problems: 3, 16, 24
    \item \textbf{C4: Biconditional Proofs (Forward and Converse)} -- Problems: 4, 19, 30
    \item \textbf{C5: Proof by Contradiction (Inequalities, Dense Sets, Non-existence)} -- Problems: 5, 9, 17, 21, 28
    \item \textbf{C6: Set Theory \& Logic (Subsets, Partitions, De Morgan's)} -- Problems: 22
    \item \textbf{C7: Diophantine \& Modular Arithmetic} -- Problems: 1, 4, 5, 16, 17, 27, 30
    \item \textbf{C8: Inequalities (AM-GM, Completing the Square)} -- Problems: 8, 23
    \item \textbf{C9: Abstract Modeling: Quant Finance (Arbitrage, Risk-Neutral Measure)} -- Problems: 10, 11, 12
    \item \textbf{C10: Abstract Modeling: Machine Learning (Boolean Hypercubes, Combinatorics)} -- Problems: 29
    \item \textbf{C11: Diagnostic / Find the Flaw (Variable Dependency, False WLOG, Missing Converse, Div by Zero)} -- Problems: 7, 13, 14, 15, 25
    \item \textbf{C12: Counterexamples (Disproving False Statements)} -- Problems: 18
\end{itemize}

\newpage

\section{Problem Sets}

\subsection{Problem 1}
Prove that if $a$ and $b$ are odd integers, then $a^2 + b^2$ is an even integer, but it is not divisible by $4$.

\subsection{Problem 2}
Prove by contrapositive: Let $x, y \in \mathbb{Z}$. If $3x + 7y$ is even, then $x$ and $y$ are of the same parity.

\subsection{Problem 3}
Prove by exhaustive cases that for any integer $n$, the expression $n^2 + n$ is even.

\subsection{Problem 4}
Prove that for any integer $x$, $5x - 3$ is even if and only if $x$ is odd.

\subsection{Problem 5}
Prove by contradiction that there are no positive integers $x$ and $y$ such that $x^2 - y^2 = 1$.

\subsection{Problem 6}
Let $x, y \in \mathbb{Z}$. Prove that if $xy$ is even, then $x$ is even or $y$ is even.

\subsection{Problem 7}
\textbf{Find the Flaw:} Evaluate the following proof. 
\textit{Theorem:} If $x$ and $y$ are even integers, then $x+y$ is divisible by $4$.
\textit{Proof:} Assume $x$ and $y$ are even. By definition, $x = 2k$ and $y = 2k$ for some integer $k$. Then $x+y = 2k + 2k = 4k$. Since $k \in \mathbb{Z}$, $x+y$ is a multiple of $4$. $\blacksquare$

\subsection{Problem 8}
Prove that for any positive real numbers $x, y, z$, the inequality $(x+y)(y+z)(z+x) \ge 8xyz$ holds.

\subsection{Problem 9}
Prove by contradiction that the product of a non-zero rational number and an irrational number is irrational.

\subsection{Problem 10}
\textbf{Quant Finance Modeling:} Consider a one-period Binomial Options Model where the initial stock price is $S_0 = 100$. In the next period, the stock can either go up by a factor of $u = 1.2$ or down by a factor of $d = 0.8$. The risk-free interest rate is $r = 0.05$. Derive the unique risk-neutral probability measure $\mathbb{Q}(up) = q$.

\subsection{Problem 11}
\textbf{Quant Finance Modeling:} Let $S_0$ be the initial stock price, $u$ the up-factor, $d$ the down-factor, and $r$ the risk-free rate. Prove by constructing a specific portfolio that if $1+r \le d$, an arbitrage opportunity exists.

\subsection{Problem 12}
\textbf{Quant Finance Modeling:} Using the same parameters as Problem 11, prove that if $1+r \ge u$, an arbitrage opportunity exists.

\subsection{Problem 13}
\textbf{Find the Flaw:} Evaluate the following proof.
\textit{Theorem:} For all $x \in \mathbb{R}$, if $x^2 > x$, then $x > 1$.
\textit{Proof:} Assume $x^2 > x$. Dividing both sides of the inequality by $x$ yields $x > 1$. Thus, the theorem holds. $\blacksquare$

\subsection{Problem 14}
\textbf{Find the Flaw:} Evaluate the following proof.
\textit{Theorem:} Let $A$ and $B$ be sets. $A = B \iff A \subseteq B \land B \subseteq A$.
\textit{Proof:} Assume $A = B$. Let $x \in A$. Since $A = B$, $x \in B$, which implies $A \subseteq B$. Similarly, let $y \in B$. Since $A = B$, $y \in A$, which implies $B \subseteq A$. Therefore, the equivalence holds. $\blacksquare$

\subsection{Problem 15}
\textbf{Find the Flaw:} Evaluate the following proof.
\textit{Theorem:} Let $x, y \in \mathbb{Z}$. If $x+y$ is even, then both $x$ and $y$ are even.
\textit{Proof:} Assume $x+y$ is even. Without loss of generality (WLOG), assume $x$ is even. Since $x$ is even, $x = 2k$. For $x+y$ to be even, $y$ must also be even ($y = 2m$), because the sum of an even and an odd is odd. Thus, both $x$ and $y$ are even. $\blacksquare$

\subsection{Problem 16}
Prove by cases that the square of any integer is of the form $3k$ or $3k+1$ for some integer $k$.

\subsection{Problem 17}
Prove by contradiction that if $a, b, c \in \mathbb{Z}$ such that $a^2 + b^2 = c^2$, then at least one of $a$ or $b$ must be even.

\subsection{Problem 18}
Find a counterexample to the following statement: "For all irrational numbers $x$ and $y$, the sum $x+y$ is irrational."

\subsection{Problem 19}
Prove that for any integer $x$, $x^3$ is even if and only if $x$ is even.

\subsection{Problem 20}
Prove by contrapositive: Let $x \in \mathbb{Z}$. If $x^2 - 6x + 5$ is even, then $x$ is odd.

\subsection{Problem 21}
Prove by contradiction that $\sqrt{3}$ is irrational.

\subsection{Problem 22}
Let $A, B, C$ be sets. Prove that $A \setminus (B \cup C) = (A \setminus B) \cap (A \setminus C)$.

\subsection{Problem 23}
Prove that for any real numbers $x$ and $y$, $x^2 + y^2 \ge 2xy$.

\subsection{Problem 24}
Prove by cases that for any real numbers $x$ and $y$, $|x+y| \le |x| + |y|$.

\subsection{Problem 25}
\textbf{Find the Flaw:} Evaluate the following proof.
\textit{Theorem:} $2 = 1$.
\textit{Proof:} Let $x = 1$. Then $x^2 = x$. Subtracting $1$ from both sides gives $x^2 - 1 = x - 1$. Factoring the left side yields $(x-1)(x+1) = x-1$. Dividing both sides by $(x-1)$ leaves $x+1 = 1$. Substituting $x=1$ gives $2 = 1$. $\blacksquare$

\subsection{Problem 26}
Prove by contrapositive: Let $x, y \in \mathbb{Z}$. If $x+y > 20$, then $x > 10$ or $y > 10$.

\subsection{Problem 27}
Prove that if $n$ is an odd integer, then $8$ divides $n^2 - 1$.

\subsection{Problem 28}
Prove by contradiction that there is no smallest positive rational number.

\subsection{Problem 29}
\textbf{Machine Learning / Boolean Hypercubes:} Let $S = \{0, 1\}^3$ be the state space of a 3-bit boolean hypercube. Prove that any subset of $5$ elements from $S$ must contain at least one pair of elements that differ in exactly one bit.

\subsection{Problem 30}
Let $x, y \in \mathbb{Z}$. Prove that $x^2 + y^2 \equiv 0 \pmod 4$ if and only if $x$ and $y$ are both even.

\newpage

\section{Detailed Solutions}

\subsection{Solution 1}
\begin{proof}
We proceed by direct proof. Assume $a$ and $b$ are odd integers. By definition, there exist integers $k$ and $m$ such that $a = 2k + 1$ and $b = 2m + 1$. We evaluate $a^2 + b^2$:
\begin{align*}
a^2 + b^2 &= (2k + 1)^2 + (2m + 1)^2 \\
&= (4k^2 + 4k + 1) + (4m^2 + 4m + 1) \\
&= 4k^2 + 4k + 4m^2 + 4m + 2 \\
&= 2(2k^2 + 2k + 2m^2 + 2m + 1).
\end{align*}
Since $2k^2 + 2k + 2m^2 + 2m + 1 \in \mathbb{Z}$, $a^2 + b^2$ is of the form $2p$ for an integer $p$, meaning it is even. 

To show it is not divisible by $4$, we factor out $4$ from the first four terms:
\[ a^2 + b^2 = 4(k^2 + k + m^2 + m) + 2. \]
Let $q = k^2 + k + m^2 + m$, where $q \in \mathbb{Z}$. Then $a^2 + b^2 = 4q + 2$. When divided by $4$, the remainder is $2$ (not $0$). Therefore, $a^2 + b^2$ is not divisible by $4$.
\end{proof}

\subsection{Solution 2}
\begin{proof}
We proceed by contrapositive. The statement is $P \implies Q$, where $P$ is "$3x + 7y$ is even" and $Q$ is "$x$ and $y$ are of the same parity". The contrapositive is $\neg Q \implies \neg P$, which translates to: "If $x$ and $y$ are of opposite parity, then $3x + 7y$ is odd."

Assume $x$ and $y$ are of opposite parity. We consider two cases:
\textbf{Case 1:} $x$ is even and $y$ is odd. 
Let $x = 2k$ and $y = 2m + 1$ for $k, m \in \mathbb{Z}$.
\begin{align*}
3x + 7y &= 3(2k) + 7(2m + 1) \\
&= 6k + 14m + 7 \\
&= 6k + 14m + 6 + 1 \\
&= 2(3k + 7m + 3) + 1.
\end{align*}
Since $3k + 7m + 3 \in \mathbb{Z}$, $3x + 7y$ is odd.

\textbf{Case 2:} $x$ is odd and $y$ is even.
Let $x = 2k + 1$ and $y = 2m$ for $k, m \in \mathbb{Z}$.
\begin{align*}
3x + 7y &= 3(2k + 1) + 7(2m) \\
&= 6k + 3 + 14m \\
&= 6k + 14m + 2 + 1 \\
&= 2(3k + 7m + 1) + 1.
\end{align*}
Since $3k + 7m + 1 \in \mathbb{Z}$, $3x + 7y$ is odd.

In all cases, $3x + 7y$ is odd. The contrapositive is proven, so the original statement holds.
\end{proof}

\subsection{Solution 3}
\begin{proof}
We proceed by exhaustive cases on the parity of $n$. The domain $\mathbb{Z}$ is partitioned into even and odd integers.
\textbf{Case 1:} $n$ is even. Let $n = 2k$ for $k \in \mathbb{Z}$.
\begin{align*}
n^2 + n &= (2k)^2 + (2k) \\
&= 4k^2 + 2k \\
&= 2(2k^2 + k).
\end{align*}
Since $2k^2 + k \in \mathbb{Z}$, $n^2 + n$ is even.

\textbf{Case 2:} $n$ is odd. Let $n = 2k + 1$ for $k \in \mathbb{Z}$.
\begin{align*}
n^2 + n &= (2k + 1)^2 + (2k + 1) \\
&= (4k^2 + 4k + 1) + (2k + 1) \\
&= 4k^2 + 6k + 2 \\
&= 2(2k^2 + 3k + 1).
\end{align*}
Since $2k^2 + 3k + 1 \in \mathbb{Z}$, $n^2 + n$ is even.
Since $n^2 + n$ is even in all possible cases, the theorem holds for all $n \in \mathbb{Z}$.
\end{proof}

\subsection{Solution 4}
\begin{proof}
We must prove both directions of the biconditional ($P \iff Q$).

\textbf{Forward Direction ($\implies$):} If $5x - 3$ is even, then $x$ is odd.
We prove this by contrapositive. Assume $x$ is even. Then $x = 2k$ for $k \in \mathbb{Z}$.
\begin{align*}
5x - 3 &= 5(2k) - 3 \\
&= 10k - 3 \\
&= 10k - 4 + 1 \\
&= 2(5k - 2) + 1.
\end{align*}
Since $5k - 2 \in \mathbb{Z}$, $5x - 3$ is odd. The contrapositive holds.

\textbf{Backward Direction ($\impliedby$):} If $x$ is odd, then $5x - 3$ is even.
We proceed by direct proof. Assume $x$ is odd. Then $x = 2k + 1$ for $k \in \mathbb{Z}$.
\begin{align*}
5x - 3 &= 5(2k + 1) - 3 \\
&= 10k + 5 - 3 \\
&= 10k + 2 \\
&= 2(5k + 1).
\end{align*}
Since $5k + 1 \in \mathbb{Z}$, $5x - 3$ is even.
Both directions are proven.
\end{proof}

\subsection{Solution 5}
\begin{proof}
We proceed by contradiction. Assume there exist positive integers $x, y \in \mathbb{Z}^+$ such that $x^2 - y^2 = 1$.
We can factor the left side as a difference of squares:
\[ (x - y)(x + y) = 1. \]
Since $x$ and $y$ are integers, both $(x - y)$ and $(x + y)$ must be integers. The only integer factors of $1$ are $(1)(1)$ and $(-1)(-1)$. 
Because $x \ge 1$ and $y \ge 1$, their sum must be $x + y \ge 2$. 
However, for the product to be $1$, we must have $x + y = 1$ or $x + y = -1$. 
This implies $2 \le x + y = 1$, which is a strict contradiction ($2 \le 1$ is false). 
Therefore, no such positive integers $x$ and $y$ exist.
\end{proof}

\subsection{Solution 6}
\begin{proof}
We proceed by contrapositive. The statement is $P \implies (Q \lor R)$. The contrapositive is $\neg(Q \lor R) \implies \neg P$. By De Morgan's Laws, $\neg(Q \lor R) \equiv \neg Q \land \neg R$. 
Thus, we must prove: "If $x$ is odd and $y$ is odd, then $xy$ is odd."

Assume $x$ and $y$ are both odd. Then $x = 2k + 1$ and $y = 2m + 1$ for $k, m \in \mathbb{Z}$.
\begin{align*}
xy &= (2k + 1)(2m + 1) \\
&= 4km + 2k + 2m + 1 \\
&= 2(2km + k + m) + 1.
\end{align*}
Since $2km + k + m \in \mathbb{Z}$, $xy$ is odd. The contrapositive is proven.
\end{proof}

\subsection{Solution 7}
\textbf{Solution Analysis:} The flaw is a \textbf{Variable Dependency Error}. By assigning $x = 2k$ and $y = 2k$, the author implicitly assumed that $x = y$. Two independent variables must be assigned independent arbitrary constants (e.g., $x = 2k$ and $y = 2m$). The correct sum would be $2(k+m)$, which proves it is even, but not necessarily a multiple of $4$ (e.g., $x=2, y=4 \implies x+y=6$, which is not divisible by $4$). The theorem itself is false.

\subsection{Solution 8}
\begin{proof}
We proceed by applying the AM-GM (Arithmetic Mean-Geometric Mean) inequality. For any positive real numbers $a, b$, the AM-GM inequality states that $\frac{a+b}{2} \ge \sqrt{ab}$, which implies $a+b \ge 2\sqrt{ab}$.
Applying this to the pairs $(x,y)$, $(y,z)$, and $(z,x)$, we get:
\begin{align*}
x + y &\ge 2\sqrt{xy} \\
y + z &\ge 2\sqrt{yz} \\
z + x &\ge 2\sqrt{zx}
\end{align*}
Since $x, y, z$ are positive, all terms are positive. We can multiply the three inequalities together:
\[ (x+y)(y+z)(z+x) \ge (2\sqrt{xy})(2\sqrt{yz})(2\sqrt{zx}) \]
\[ (x+y)(y+z)(z+x) \ge 8\sqrt{x^2 y^2 z^2} \]
Since $x, y, z > 0$, $\sqrt{x^2 y^2 z^2} = xyz$. Thus:
\[ (x+y)(y+z)(z+x) \ge 8xyz. \]
\end{proof}

\subsection{Solution 9}
\begin{proof}
We proceed by contradiction. Let $a \in \mathbb{Q}$ with $a \neq 0$, and let $b \notin \mathbb{Q}$. Assume that their product $ab$ is rational, meaning $ab \in \mathbb{Q}$.
Since $a \in \mathbb{Q}$ and $a \neq 0$, $a = \frac{p}{q}$ for integers $p, q$ with $p \neq 0, q \neq 0$.
Since $ab \in \mathbb{Q}$, $ab = \frac{m}{n}$ for integers $m, n$ with $n \neq 0$.
We can solve for $b$:
\[ b = \frac{ab}{a} = \frac{\frac{m}{n}}{\frac{p}{q}} = \frac{mq}{np} \]
Since $m, q, n, p$ are integers, and $n \neq 0, p \neq 0 \implies np \neq 0$, the product $mq$ is an integer and $np$ is a non-zero integer. 
This implies that $b$ can be expressed as a ratio of integers, meaning $b \in \mathbb{Q}$. 
This contradicts the initial premise that $b \notin \mathbb{Q}$. Therefore, the product $ab$ must be irrational.
\end{proof}

\subsection{Solution 10}
\begin{proof}
In a risk-neutral measure $\mathbb{Q}$, the expected discounted payoff of the stock must equal its current price. Thus, $\mathbb{E}^{\mathbb{Q}}\left[\frac{S_1}{1+r}\right] = S_0$.
This gives the equation:
\[ \frac{q(uS_0) + (1-q)(dS_0)}{1+r} = S_0 \]
Divide both sides by $S_0$ (since $S_0 \neq 0$):
\[ q u + (1-q) d = 1+r \]
Substitute the given values ($u = 1.2, d = 0.8, r = 0.05$):
\[ q(1.2) + (1-q)(0.8) = 1.05 \]
\[ 1.2q + 0.8 - 0.8q = 1.05 \]
\[ 0.4q = 0.25 \]
\[ q = \frac{0.25}{0.40} = \frac{5}{8} = 0.625 \]
The risk-neutral probability measure is $q = 0.625$.
\end{proof}

\subsection{Solution 11}
\begin{proof}
We proceed by constructing an arbitrage portfolio. Assume $1+r \le d$. 
At $t=0$, borrow $S_0$ dollars at the risk-free rate $r$ and purchase exactly $1$ share of the stock. The initial cash flow is $S_0 - S_0 = 0$.
At $t=1$, we must repay the loan, which costs $S_0(1+r)$. We also hold $1$ share of stock, which is worth $S_1 \in \{uS_0, dS_0\}$.
We evaluate the portfolio payoff in all states of the world (Proof by Cases):
\textbf{Case 1 (Down State):} The stock is worth $dS_0$. The net profit is $dS_0 - S_0(1+r) = S_0(d - (1+r))$. Since $1+r \le d$, $d - (1+r) \ge 0$. The profit is $\ge 0$.
\textbf{Case 2 (Up State):} The stock is worth $uS_0$. The net profit is $uS_0 - S_0(1+r) = S_0(u - (1+r))$. By definition of the model, $u > d$, so $u > 1+r$. The profit is strictly $> 0$.
Since the portfolio requires zero initial investment, never loses money, and has a strictly positive probability of making money, it constitutes an arbitrage opportunity.
\end{proof}

\subsection{Solution 12}
\begin{proof}
We proceed by constructing an arbitrage portfolio. Assume $1+r \ge u$.
At $t=0$, short-sell $1$ share of the stock to receive $S_0$, and lend this $S_0$ at the risk-free rate $r$. The initial cash flow is $S_0 - S_0 = 0$.
At $t=1$, we receive the loan repayment of $S_0(1+r)$. We must also buy back the stock to cover the short position, which costs $S_1 \in \{uS_0, dS_0\}$.
We evaluate the portfolio payoff in all states of the world:
\textbf{Case 1 (Up State):} The stock is worth $uS_0$. The net profit is $S_0(1+r) - uS_0 = S_0((1+r) - u)$. Since $1+r \ge u$, $(1+r) - u \ge 0$. The profit is $\ge 0$.
\textbf{Case 2 (Down State):} The stock is worth $dS_0$. The net profit is $S_0(1+r) - dS_0 = S_0((1+r) - d)$. By definition, $u > d$, so $1+r > d$. The profit is strictly $> 0$.
Since the portfolio requires zero initial investment, never loses money, and has a strictly positive probability of making money, it constitutes an arbitrage opportunity.
\end{proof}

\subsection{Solution 13}
\textbf{Solution Analysis:} The flaw is an \textbf{Algebraic Manipulation Error (Sign Flip)}. The proof divides the inequality $x^2 > x$ by $x$. In inequalities, dividing by a variable is invalid unless the sign of the variable is known, because dividing by a negative number flips the inequality sign. If $x < 0$, dividing by $x$ yields $x < 1$. A counterexample to the theorem is $x = -2$. $(-2)^2 = 4 > -2$, but $-2 \not> 1$.

\subsection{Solution 14}
\textbf{Solution Analysis:} The flaw is a \textbf{Missing Converse}. The proof only establishes the forward direction ($\implies$): assuming $A = B$ and deriving $A \subseteq B \land B \subseteq A$. Proving a biconditional ($\iff$) requires a separate logical block assuming $A \subseteq B \land B \subseteq A$ and deriving $A = B$.

\subsection{Solution 15}
\textbf{Solution Analysis:} The flaw is a \textbf{False WLOG (Without Loss of Generality) / Converse Error}. The author assumes that because $x+y$ is even, one of the variables must be even. This is logically invalid; $x$ and $y$ could both be odd (e.g., $3+5 = 8$). WLOG can only be used when variables are symmetric in the premise \textit{and} the assumption doesn't restrict the universal domain of the premise. The theorem itself is false.

\subsection{Solution 16}
\begin{proof}
We proceed by cases, partitioning the integers modulo $3$. For any integer $x$, $x$ can be expressed as $3q$, $3q+1$, or $3q+2$ for some integer $q$.
\textbf{Case 1:} $x = 3q$.
\[ x^2 = (3q)^2 = 9q^2 = 3(3q^2). \]
Let $k = 3q^2$. Then $x^2 = 3k$.

\textbf{Case 2:} $x = 3q+1$.
\[ x^2 = (3q+1)^2 = 9q^2 + 6q + 1 = 3(3q^2 + 2q) + 1. \]
Let $k = 3q^2 + 2q$. Then $x^2 = 3k + 1$.

\textbf{Case 3:} $x = 3q+2$.
\[ x^2 = (3q+2)^2 = 9q^2 + 12q + 4 = 9q^2 + 12q + 3 + 1 = 3(3q^2 + 4q + 1) + 1. \]
Let $k = 3q^2 + 4q + 1$. Then $x^2 = 3k + 1$.

In all cases, $x^2$ is of the form $3k$ or $3k+1$.
\end{proof}

\subsection{Solution 17}
\begin{proof}
We proceed by contradiction. Assume there exist integers $a, b, c$ such that $a^2 + b^2 = c^2$, but neither $a$ nor $b$ is even. This means both $a$ and $b$ are odd.
By definition, $a = 2k+1$ and $b = 2m+1$.
\begin{align*}
a^2 + b^2 &= (2k+1)^2 + (2m+1)^2 \\
&= 4k^2 + 4k + 1 + 4m^2 + 4m + 1 \\
&= 4(k^2 + k + m^2 + m) + 2.
\end{align*}
This shows that $a^2 + b^2 \equiv 2 \pmod 4$.
Now consider $c^2$. Since $c$ is an integer, it is either even or odd.
If $c$ is even ($c=2p$), $c^2 = 4p^2 \equiv 0 \pmod 4$.
If $c$ is odd ($c=2p+1$), $c^2 = 4p^2+4p+1 = 4(p^2+p)+1 \equiv 1 \pmod 4$.
Thus, for any integer $c$, $c^2$ is congruent to $0$ or $1$ modulo $4$.
However, we established $a^2 + b^2 \equiv 2 \pmod 4$. Therefore, $a^2 + b^2 \neq c^2$, which contradicts our initial assumption. Thus, at least one of $a$ or $b$ must be even.
\end{proof}

\subsection{Solution 18}
\textbf{Counterexample:} Let $x = \sqrt{2}$ and $y = -\sqrt{2}$. Both $x$ and $y$ are irrational numbers. However, their sum is $x + y = \sqrt{2} + (-\sqrt{2}) = 0$. Since $0$ can be expressed as $\frac{0}{1}$, it is a rational number. This disproves the statement.

\subsection{Solution 19}
\begin{proof}
We must prove both directions of the biconditional.
\textbf{Forward ($\implies$):} If $x^3$ is even, then $x$ is even.
We use the contrapositive: If $x$ is odd, then $x^3$ is odd.
Assume $x$ is odd. $x = 2k+1$.
\begin{align*}
x^3 &= (2k+1)^3 \\
&= 8k^3 + 12k^2 + 6k + 1 \\
&= 2(4k^3 + 6k^2 + 3k) + 1.
\end{align*}
Since $4k^3 + 6k^2 + 3k \in \mathbb{Z}$, $x^3$ is odd. The contrapositive holds.

\textbf{Backward ($\impliedby$):} If $x$ is even, then $x^3$ is even.
We use direct proof. Assume $x$ is even. $x = 2k$.
\[ x^3 = (2k)^3 = 8k^3 = 2(4k^3). \]
Since $4k^3 \in \mathbb{Z}$, $x^3$ is even. Both directions hold.
\end{proof}

\subsection{Solution 20}
\begin{proof}
We proceed by contrapositive. The statement is "If $x^2 - 6x + 5$ is even, then $x$ is odd." The contrapositive is "If $x$ is even, then $x^2 - 6x + 5$ is odd."
Assume $x$ is even. Then $x = 2k$ for some integer $k$.
\begin{align*}
x^2 - 6x + 5 &= (2k)^2 - 6(2k) + 5 \\
&= 4k^2 - 12k + 5 \\
&= 4k^2 - 12k + 4 + 1 \\
&= 2(2k^2 - 6k + 2) + 1.
\end{align*}
Since $2k^2 - 6k + 2 \in \mathbb{Z}$, the expression is of the form $2m + 1$, which means $x^2 - 6x + 5$ is odd. The contrapositive is proven.
\end{proof}

\subsection{Solution 21}
\begin{proof}
We proceed by contradiction. Assume $\sqrt{3}$ is rational. Then it can be expressed as an irreducible fraction $\sqrt{3} = \frac{p}{q}$, where $p, q \in \mathbb{Z}$, $q \neq 0$, and $\gcd(p,q) = 1$ (they share no common factors).
Squaring both sides yields:
\[ 3 = \frac{p^2}{q^2} \implies p^2 = 3q^2. \]
This implies that $p^2$ is a multiple of $3$. Since $3$ is prime, if $3$ divides $p^2$, then $3$ must divide $p$. Thus, $p = 3k$ for some integer $k$.
Substitute $p = 3k$ back into the equation:
\[ (3k)^2 = 3q^2 \implies 9k^2 = 3q^2 \implies 3k^2 = q^2. \]
This implies that $q^2$ is a multiple of $3$. By the same logic, $3$ must divide $q$.
If $3$ divides both $p$ and $q$, then $\gcd(p,q) \ge 3$. This contradicts our initial assumption that $\gcd(p,q) = 1$. Therefore, $\sqrt{3}$ must be irrational.
\end{proof}

\subsection{Solution 22}
\begin{proof}
We prove the set equality by establishing logical equivalence for an arbitrary element $x$.
\begin{align*}
x \in A \setminus (B \cup C) &\iff x \in A \land x \notin (B \cup C) \quad \text{(Definition of set difference)} \\
&\iff x \in A \land \neg(x \in B \lor x \in C) \quad \text{(Definition of union)} \\
&\iff x \in A \land (x \notin B \land x \notin C) \quad \text{(De Morgan's Law)} \\
&\iff (x \in A \land x \notin B) \land (x \in A \land x \notin C) \quad \text{(Idempotence and Associativity of $\land$)} \\
&\iff x \in (A \setminus B) \land x \in (A \setminus C) \quad \text{(Definition of set difference)} \\
&\iff x \in (A \setminus B) \cap (A \setminus C) \quad \text{(Definition of intersection)}
\end{align*}
Since the logical equivalence holds for all $x$, the sets are equal: $A \setminus (B \cup C) = (A \setminus B) \cap (A \setminus C)$.
\end{proof}

\subsection{Solution 23}
\begin{proof}
We proceed by direct proof using the property of squares. For any real numbers $x$ and $y$, the square of their difference must be non-negative:
\[ (x - y)^2 \ge 0 \]
Expanding the polynomial yields:
\[ x^2 - 2xy + y^2 \ge 0 \]
Adding $2xy$ to both sides gives:
\[ x^2 + y^2 \ge 2xy \]
This completes the proof.
\end{proof}

\subsection{Solution 24}
\begin{proof}
We proceed by cases based on the sign of $(x+y)$.
\textbf{Case 1:} $x+y \ge 0$.
Then $|x+y| = x+y$.
By definition of absolute value, $x \le |x|$ and $y \le |y|$.
Adding these inequalities yields $x+y \le |x| + |y|$.
Thus, $|x+y| \le |x| + |y|$.

\textbf{Case 2:} $x+y < 0$.
Then $|x+y| = -(x+y) = -x - y$.
By definition of absolute value, $-x \le |x|$ and $-y \le |y|$.
Adding these inequalities yields $-x - y \le |x| + |y|$.
Thus, $|x+y| \le |x| + |y|$.

In all cases, the Triangle Inequality $|x+y| \le |x| + |y|$ holds.
\end{proof}

\subsection{Solution 25}
\textbf{Solution Analysis:} The flaw is \textbf{Division by Zero}. The proof assumes $x = 1$. In the step where both sides are divided by $(x-1)$, the author is dividing by $(1-1) = 0$. Division by zero is undefined, rendering all subsequent logical deductions invalid.

\subsection{Solution 26}
\begin{proof}
We proceed by contrapositive. The statement is $P \implies (Q \lor R)$. The contrapositive is $\neg(Q \lor R) \implies \neg P$. By De Morgan's Laws, this is $(\neg Q \land \neg R) \implies \neg P$.
We must prove: "If $x \le 10$ and $y \le 10$, then $x+y \le 20$."
Assume $x \le 10$ and $y \le 10$.
Adding the two inequalities yields:
\[ x + y \le 10 + 10 \implies x + y \le 20. \]
The contrapositive is proven, so the original statement holds.
\end{proof}

\subsection{Solution 27}
\begin{proof}
We proceed by direct proof. Assume $n$ is an odd integer. Then $n = 2k+1$ for some integer $k$.
We evaluate $n^2 - 1$:
\begin{align*}
n^2 - 1 &= (2k+1)^2 - 1 \\
&= 4k^2 + 4k + 1 - 1 \\
&= 4k^2 + 4k \\
&= 4k(k+1).
\end{align*}
Notice that $k$ and $(k+1)$ are consecutive integers. Exactly one of them must be even. Therefore, their product $k(k+1)$ is an even integer. Let $k(k+1) = 2m$ for some integer $m$.
Substitute this back into the equation:
\[ n^2 - 1 = 4(2m) = 8m. \]
Since $m \in \mathbb{Z}$, $n^2 - 1$ is a multiple of $8$, meaning $8$ divides $n^2 - 1$.
\end{proof}

\subsection{Solution 28}
\begin{proof}
We proceed by contradiction. Assume there exists a smallest positive rational number, let's call it $r$. Since $r$ is positive and rational, $r \in \mathbb{Q}$ and $r > 0$.
Consider the number $x = \frac{r}{2}$. 
Since $r$ is rational, $r = \frac{p}{q}$ for integers $p, q$. Then $x = \frac{p}{2q}$, which is a ratio of integers, so $x \in \mathbb{Q}$.
Since $r > 0$, dividing by $2$ yields $x > 0$.
Furthermore, since $2 > 1$, dividing $r$ by $2$ yields $x < r$.
Thus, $x$ is a positive rational number that is strictly smaller than $r$. This contradicts the assumption that $r$ is the smallest positive rational number. Therefore, no such number exists.
\end{proof}

\subsection{Solution 29}
\begin{proof}
We proceed by partitioning the state space and using the Pigeonhole Principle. The set $S = \{0,1\}^3$ contains $2^3 = 8$ elements.
We can partition $S$ into $4$ disjoint subsets (pairs), where the elements in each pair differ by exactly the last bit:
\begin{align*}
P_1 &= \{000, 001\} \\
P_2 &= \{010, 011\} \\
P_3 &= \{100, 101\} \\
P_4 &= \{110, 111\}
\end{align*}
Notice that $P_1 \cup P_2 \cup P_3 \cup P_4 = S$ and $P_i \cap P_j = \emptyset$ for $i \neq j$.
If we select any $5$ elements from $S$, we are distributing $5$ items into $4$ available sets (the partitions $P_1, P_2, P_3, P_4$).
By the Pigeonhole Principle, at least one set must contain $\lceil 5/4 \rceil = 2$ elements.
Since each set $P_i$ contains exactly two elements that differ in exactly one bit, selecting both elements from that set guarantees we have a pair that differs by exactly one bit.
\end{proof}

\subsection{Solution 30}
\begin{proof}
We must prove both directions of the biconditional.
\textbf{Backward ($\impliedby$):} If $x$ and $y$ are both even, then $x^2 + y^2 \equiv 0 \pmod 4$.
Assume $x = 2k$ and $y = 2m$.
\[ x^2 + y^2 = (2k)^2 + (2m)^2 = 4k^2 + 4m^2 = 4(k^2 + m^2). \]
Since $k^2 + m^2 \in \mathbb{Z}$, $x^2 + y^2$ is a multiple of $4$, so $x^2 + y^2 \equiv 0 \pmod 4$.

\textbf{Forward ($\implies$):} If $x^2 + y^2 \equiv 0 \pmod 4$, then $x$ and $y$ are both even.
We proceed by contrapositive: If it is not the case that both are even, then $x^2 + y^2 \not\equiv 0 \pmod 4$.
Assume it is not the case that both are even. We have two cases:
\textbf{Case 1:} Both $x$ and $y$ are odd. $x = 2k+1, y = 2m+1$.
\begin{align*}
x^2 + y^2 &= (2k+1)^2 + (2m+1)^2 \\
&= 4k^2 + 4k + 1 + 4m^2 + 4m + 1 \\
&= 4(k^2 + k + m^2 + m) + 2.
\end{align*}
This means $x^2 + y^2 \equiv 2 \pmod 4$, which is not $0 \pmod 4$.

\textbf{Case 2:} One is even, one is odd. WLOG, $x = 2k, y = 2m+1$.
\begin{align*}
x^2 + y^2 &= (2k)^2 + (2m+1)^2 \\
&= 4k^2 + 4m^2 + 4m + 1 \\
&= 4(k^2 + m^2 + m) + 1.
\end{align*}
This means $x^2 + y^2 \equiv 1 \pmod 4$, which is not $0 \pmod 4$.
In all cases, the contrapositive holds. Both directions are proven.
\end{proof}

\newpage

\section{Concept Checklist \& Question Mapping (Reference)}
\label{sec:checklist_end}

\begin{itemize}
    \item \textbf{C1: Direct Proof (Parity \& Algebra)} -- Problems: 1, 27
    \item \textbf{C2: Proof by Contrapositive (Implications \& De Morgan's)} -- Problems: 2, 6, 20, 26
    \item \textbf{C3: Proof by Cases (Exhaustion \& Partitions)} -- Problems: 3, 16, 24
    \item \textbf{C4: Biconditional Proofs (Forward and Converse)} -- Problems: 4, 19, 30
    \item \textbf{C5: Proof by Contradiction (Inequalities, Dense Sets, Non-existence)} -- Problems: 5, 9, 17, 21, 28
    \item \textbf{C6: Set Theory \& Logic (Subsets, Partitions, De Morgan's)} -- Problems: 22
    \item \textbf{C7: Diophantine \& Modular Arithmetic} -- Problems: 1, 4, 5, 16, 17, 27, 30
    \item \textbf{C8: Inequalities (AM-GM, Completing the Square)} -- Problems: 8, 23
    \item \textbf{C9: Abstract Modeling: Quant Finance (Arbitrage, Risk-Neutral Measure)} -- Problems: 10, 11, 12
    \item \textbf{C10: Abstract Modeling: Machine Learning (Boolean Hypercubes, Combinatorics)} -- Problems: 29
    \item \textbf{C11: Diagnostic / Find the Flaw (Variable Dependency, False WLOG, Missing Converse, Div by Zero)} -- Problems: 7, 13, 14, 15, 25
    \item \textbf{C12: Counterexamples (Disproving False Statements)} -- Problems: 18
\end{itemize}

\end{document}